%% Question bank chapter file - generated by bank_tools.py build (do not edit while the bank is still being built:
%% it is regenerated from the per-test data in Question Bank\data). Every question is one bankq block:
%%   \begin{bankq}[id=..., source={...}, paper=.., qno=.., type=..., subtopic={...}, answer={...}, answer_src=key|solved,
%%                 working={...}, difficulty=1-10, diff_parts={K S I C F}, diff_note={...},
%%                 group=<passage id or empty>, also_in={duplicate copies in other tests}]
%% difficulty = min(10, K+S+I+C+F-1): K concepts 1-3, S steps 1-4, I insight 0-2, C computation 0-2, F format 0-1
%%   ...question LaTeX, options, figures (\bfig / \bfigi -> img/)...
%%   \end{bankq}
%% Passages shared by several questions are bankpassage blocks; their questions carry group=<passage id>.
\documentclass[10pt]{article}
\usepackage[a4paper,margin=1.8cm]{geometry}
\usepackage[T1]{fontenc}
\usepackage{amsmath,amssymb,mathtools}
\usepackage{graphicx}
\usepackage[export]{adjustbox}
\usepackage{xcolor,needspace,keyval}
\usepackage{enumitem}
\usepackage{booktabs,array,tabularx}
\usepackage[hidelinks]{hyperref}
\setlength{\parindent}{0pt}
\setlength{\parskip}{3pt}
\makeatletter
\define@key{bq}{id}{\def\bq@id{#1}}\define@key{bq}{source}{\def\bq@src{#1}}\define@key{bq}{paper}{\def\bq@paper{#1}}
\define@key{bq}{qno}{\def\bq@qno{#1}}\define@key{bq}{type}{\def\bq@type{#1}}\define@key{bq}{subtopic}{\def\bq@sub{#1}}
\define@key{bq}{answer}{\def\bq@ans{#1}}\define@key{bq}{answer_src}{\def\bq@asrc{#1}}\define@key{bq}{working}{\def\bq@work{#1}}
\define@key{bq}{group}{\def\bq@group{#1}}\define@key{bq}{also_in}{\def\bq@also{#1}}
\define@key{bq}{difficulty}{\def\bq@diff{#1}}\define@key{bq}{diff_parts}{\def\bq@dparts{#1}}\define@key{bq}{diff_note}{\def\bq@dnote{#1}}
\newenvironment{bankq}[1][]{\def\bq@work{}\def\bq@also{}\def\bq@group{}\def\bq@diff{?}\def\bq@dparts{}\def\bq@dnote{}\setkeys{bq}{#1}%
  \par\needspace{4\baselineskip}\medskip\noindent{\small\textbf{\bq@id}\hfill\color{gray}[\bq@src, Paper \bq@paper, Q\bq@qno]
  \bq@type\ $\cdot$ \bq@sub\ $\cdot$ \textbf{D\bq@diff}}\par\nopagebreak\relax}
  {\par{\small\color{teal}Answer: \textbf{\bq@ans}\ifx\bq@work\@empty\else\ (solved: \bq@work)\fi
  \ifx\bq@also\@empty\else\par\color{gray}Also in: \bq@also\fi}\par}
\define@key{bp}{id}{\def\bp@id{#1}}\define@key{bp}{source}{\def\bp@src{#1}}\define@key{bp}{questions}{\def\bp@qs{#1}}
\newenvironment{bankpassage}[1][]{\setkeys{bp}{#1}\par\needspace{4\baselineskip}\bigskip\noindent
  {\small\textbf{Passage \bp@id}\hfill\color{gray}[\bp@src; questions \bp@qs]}\par\nopagebreak\relax\itshape}{\par}
\makeatother
\newcommand{\bfig}[2]{\par\smallskip{\centering\includegraphics[width=#1pt,max width=0.62\linewidth,max height=0.32\textheight]{img/#2}\par}\smallskip}
\newcommand{\bfigi}[2]{\includegraphics[width=#1pt,max width=\linewidth,max height=0.32\textheight,valign=c]{img/#2}}
\newcommand{\bmissing}[1]{\textit{[figure missing in the source conversion: \detokenize{#1}]}}
\newenvironment{qoptions}{\begin{itemize}[leftmargin=1.8em,itemsep=1pt,topsep=2pt,label={}]}{\end{itemize}}
\newcommand{\opt}[2]{\item[(#1)] #2}
\providecommand{\fwstart}{}\providecommand{\fwend}{}\providecommand{\qfullsplit}{}
\newcolumntype{L}[1]{>{\raggedright\arraybackslash}p{#1}}
\newcolumntype{C}[1]{>{\centering\arraybackslash}p{#1}}
\newcolumntype{R}[1]{>{\raggedleft\arraybackslash}p{#1}}
\providecommand{\sectionhdr}[1]{}\providecommand{\subsectionhdr}[1]{}


\begin{document}
{\Large\bfseries Hydrogen and its Compounds --- Question Bank}\par\smallskip
{\small Chemistry chapter C38. Questions from all converted test papers, grouped by source series. Answers: printed key where it matches the paper, otherwise solved (answer\_src=solved, with working). Exact repeats across papers appear once, with also\_in listing the other copies.}\par\medskip\hrule

\section*{TOAS papers}
\begin{bankq}[id=TOAS03-P2-CHEM-S3-Q04, source={TOAS - Test 03 (27 Feb 2022)}, paper=2, qno=4, type=numerical, subtopic={Water}, answer={2}, answer_src=key, difficulty=2, diff_parts={K1 S1 I0 C0 F1}, diff_note={concepts: temporary hardness | steps: equivalents of HCl used; ppm as CaCO3}]
50 ml of water sample containing temporary hardness only, required 0.1 ml of M/50 HCl for complete neutralisation. Calculate the temporary hardness of water in ppm.
\end{bankq}
\begin{bankq}[id=TOAS17-P2-CHEM-S1-Q15, source={TOAS - Test 17 (18 Sep 2022)}, paper=2, qno=15, type=single, subtopic={Water}, answer={C}, answer_src=key, difficulty=2, diff_parts={K1 S2 I0 C0 F0}, diff_note={concepts: heavy water | steps: neutron count; boiling point; valence electrons}]
Which of the following statement is correct for H$_2$O and D$_2$O?
\begin{qoptions}
\opt{A}{D$_2$O has higher number of valence shell electron than H$_2$O}
\opt{B}{D$_2$O has lower value of boiling point than H$_2$O}
\opt{C}{D$_2$O has higher number of neutrons than H$_2$O}
\opt{D}{H$_2$O has smaller number of protons than D$_2$O}
\end{qoptions}
\end{bankq}

\section*{TNAS papers}
\begin{bankq}[id=TNAS07-P1-CHEM-S1-Q01, source={TNAS - Test 07 (20 Aug 2023)}, paper=1, qno=1, type=single, subtopic={Water}, answer={D}, answer_src=key, difficulty=1, diff_parts={K1 S1 I0 C0 F0}, diff_note={concepts: properties of heavy water | steps: compare the properties}]
Select the property which have higher value for $\mathrm{H_2O}$ when compared to $\mathrm{D_2O}$.
\begin{qoptions}
\opt{A}{Density}
\opt{B}{Boiling point}
\opt{C}{Viscosity}
\opt{D}{Dielectric constant}
\end{qoptions}
\end{bankq}
\begin{bankq}[id=TNAS07-P1-CHEM-S1-Q02, source={TNAS - Test 07 (20 Aug 2023)}, paper=1, qno=2, type=single, subtopic={Hydrides}, answer={D}, answer_src=key, difficulty=2, diff_parts={K2 S1 I0 C0 F0}, diff_note={concepts: hydrogen as a fuel, storage | steps: energy content; storage methods}]
Select the correct statement regarding hydrogen as a fuel.
\begin{qoptions}
\opt{A}{The energy released on combustion per gram for hydrogen is less than LPG but greater than petrol.}
\opt{B}{The pollutants that will be released will be oxides of nitrogen for which there is no way of reducing it.}
\opt{C}{The cylinder of compressed hydrogen would weigh much lighter than the tank of petrol as hydrogen is lightest gas.}
\opt{D}{Expensive insulted tanks made from metal alloys like $\mathrm{NaNi_5}$ , $\mathrm{Ti{-}TiH_2}$ , $\mathrm{Mg{-}MgH_2}$ etc are in use for storage of dihydrogen in small quantities.}
\end{qoptions}
\end{bankq}
\begin{bankq}[id=TNAS07-P2-CHEM-S1-Q02, source={TNAS - Test 07 (20 Aug 2023)}, paper=2, qno=2, type=multi, subtopic={H2O2}, answer={B,C,D}, answer_src=key, difficulty=3, diff_parts={K1 S1 I1 C0 F1}, diff_note={concepts: roles of hydrogen peroxide | steps: oxidation numbers; oxidising or reducing in each reaction}]
In which of the following reaction(s) $\mathrm{H_2O_2}$ act as reducing agent?
\begin{qoptions}
\opt{A}{$\mathrm{PbS+4H_2O_2\longrightarrow PbSO_4+4H_2O}$}
\opt{B}{{\footnotesize$\mathrm{2KMnO_4+3H_2SO_4+5H_2O_2}\longrightarrow$ $\mathrm{2MnSO_4+K_2SO_4+}$ $\mathrm{8H_2O+5O_2}$}}
\opt{C}{$\mathrm{HOCl+H_2O_2\longrightarrow HCl+H_2O+O_2}$}
\opt{D}{$\mathrm{I_2+H_2O_2+2KOH\longrightarrow 2KI+2H_2O+O_2}$}
\end{qoptions}
\end{bankq}

\section*{Star Review (Other) papers}
\begin{bankpassage}[id=STAR01-P1-CHEM-S3-G1-2, source={Star Review (Other) - Test 01 (12 Sep 2022)}, questions={1-2}]
Born - Haber cycle for formation of solid binary metal hydride is given below along with energy profile diagram of $\mathrm{H_2}$ molecule.
\bfig{121.0}{STAR01_p1_chem_s3_para1_2_1.png}
where Me $\to$ Alkali metal or Alkaline earth metal
\bfig{97.0}{STAR01_p1_chem_s3_para1_2_2.png}
The following data can be useful.\par
$\left(\Delta H_{IE_1}\right)_{Na}=495$ kJ/mol\par
$\left(\Delta H_{IE_1}\right)_{Ca}=589$ kJ/mol\par
$\left(\Delta H_{EG}\right)_{Ca^{2+}}=-1145$ kJ/mol\par
$\left(\Delta H_{LE}\right)_{NaH}=808$ kJ/mol\par
$\left(\Delta H_{LE}\right)_{CaH_2}=2390$ kJ/mol\par
$\left(\Delta H_{E.G._1}\right)_H=-72$ kJ/mol
\end{bankpassage}
\begin{bankq}[id=STAR01-P1-CHEM-S3-Q01, source={Star Review (Other) - Test 01 (12 Sep 2022)}, paper=1, qno=1, type=numerical, subtopic={Hydrides}, answer={3.00}, answer_src=key, difficulty=7, diff_parts={K1 S2 I2 C2 F1}, diff_note={concepts: Born Haber cycle for ionic hydrides | steps: cycle for each hydride; sum of the steps; difference of formation enthalpies}, group=STAR01-P1-CHEM-S3-G1-2]
The value of $\left|\left(\Delta H_f\right)_{CaH_2}-\left(\Delta H_f\right)_{NaH}\right|$ if $\left\{\left(\Delta H_{sublimation}\right)_{Ca}-\left(\Delta H_{Sublimation}\right)_{Na}\right\}$ is 200 kJ/mol is \_\_\_\_\_\_.
\end{bankq}
\begin{bankq}[id=STAR01-P1-CHEM-S3-Q02, source={Star Review (Other) - Test 01 (12 Sep 2022)}, paper=1, qno=2, type=numerical, subtopic={Hydrides}, answer={641.00}, answer_src=key, difficulty=7, diff_parts={K1 S2 I2 C2 F1}, diff_note={concepts: enthalpy of formation of ions from the cycle | steps: steps of the process; ionisation energy; electron gain; combine}, group=STAR01-P1-CHEM-S3-G1-2]
If change in enthalpy for the process $\mathrm{Me(g)}+\dfrac n2\mathrm H_2(g)\to\mathrm{Me^{n+}(g)}+n\mathrm H^-(g)$ is $\Delta H_1$ then the value of $|\Delta H_1|$ when Me is "Na" is \_\_\_\_\_\_\_
\end{bankq}

\section*{Nurture (General) papers}
\begin{bankq}[id=NUR01-P2-CHEM-S1-Q01, source={Nurture (General) - Test 01 (12 Jun 2021)}, paper=2, qno=1, type=multi, subtopic={Hydrides}, answer={B}, answer_src=solved, working={HNO2 has nitrogen at +3, intermediate oxidation state}, difficulty=3, diff_parts={K1 S2 I0 C0 F1}, diff_note={concepts: statements on nitrogen and its compounds | steps: volatility; oxidation states; polarity; reactivity}]
Choose correct Statement:
\begin{qoptions}
\opt{A}{In group 15 Hydrides, $\mathrm{NH_3}$ is least volatile.}
\opt{B}{$\mathrm{HNO_2}$ can act as oxidiser as well as reducer.}
\opt{C}{$\mathrm{NH_3}$ is a non-polar molecule.}
\opt{D}{$\mathrm{N_2}$ is a reactive gas at room temperature.}
\end{qoptions}
\end{bankq}
\end{document}
